\documentclass[11pt,a4paper]{article}
\usepackage[T1]{fontenc}
\usepackage[utf8]{inputenc}
\usepackage[english]{babel}
\usepackage{amsmath,amssymb}
\usepackage[margin=2cm]{geometry}
\setlength{\parindent}{0pt}
\setlength{\emergencystretch}{3em}
\allowdisplaybreaks

\title{About \texttt{false\_maps\_by\_cases.tex}}
\author{}
\date{}
\begin{document}
\maketitle
This note explains what the test data in \texttt{false\_maps\_by\_cases.tex} is used for. That file is input for the engine in \texttt{main.c}; it is read by the test \texttt{test\_set\_theory.c} in this folder.

\begin{enumerate}
\item \textbf{Purpose.} Anti-regression for the maps defined by cases (paragraph 17s of \texttt{parser.c}, rules 18v of \texttt{main.c}). The file holds 28 statements, each a true one changed a little: a coefficient changed so that the map is no longer injective (items 2, 8, 9, 10, 13), a value missed by a map claimed surjective (items 3, 6, 14, 16), injectivity claimed for a map that is not injective (items 7, 15), a parameter case with the wrong value (item 11), and five whose hypotheses are not a definition: cases that overlap with different values (items 1 and 4: $n \ge 1$ instead of $n \ge 2$ in the first, $n \ge 0$ and $n \le 0$ in the second), cases that leave $0$ out (item 5), and a value outside the codomain ($|0| = 0 \notin \mathbb{N}$, item 12; $0^{2} = 0 \notin \mathbb{N}$ as a second component, item 20). Items 17--21 have values in pairs: surjectivity of $n \mapsto (n + 1, n)$, injectivity of $n \mapsto (n^{2}, n^{2})$, non-injectivity of $n \mapsto (n + 1, -n - 1)$, and a pair that is not a value. Items 22 and 23 are the true items 18 and 19 of \texttt{maps\_by\_cases.tex} with the two components of the pair written out exchanged: on $\mathbb{Z} \times \mathbb{N}_{0}$ and $\mathbb{N} \times \mathbb{Z}$ the two sides differ, so a reader that exchanged the components of a variable over a set of pairs would prove them (an engine broken on purpose in this way proved both, and gave every other verdict of this file right). Items 24--28 are about the iterates $f^{n}$: two wrong closed forms (the second agrees with the right one for $n = 1, 2$ and fails for $n = 3$), two maps whose values are not all in the domain, so that $f \circ f^{n}$ is not defined ($f(1) = \frac{1}{2} \notin \mathbb{Z}$; $f(1) = 0 \notin \mathbb{N}$, while the closed form $x - n$ satisfies the recursion as an identity, and only the domain stops it), and $f^{2} = f$ in place of $f^{1} = f$, which leaves $f^{1}$ undefined.

\item \textbf{The traps.} Items 1, 4, 12 and 20 have hypotheses that no map satisfies: as implications they are true for the empty reason, and a careless engine would prove them, or prove anything else from them. Item 5 leaves $f(0)$ free, so its conclusion holds for some maps and fails for others. The engine must refuse all five: before any claim it proves that the cases cover the domain, agree where they overlap and give values in the codomain, and here it cannot. Item 11 is false only for some values of the parameter: for $k = 1$ the map $n \mapsto n + 1$ is injective, for $k = 2$ it is not ($f(0) = 1 = f(1)$). The engine must not refute it for every $k$, nor prove it: it refutes it with $k = 2$, after checking that for $k = 2$ the cases define a map.

\item \textbf{What is being tested.} That none is proved. The others are refuted, with a counterexample.

\item \textbf{How to run it.} From the top folder of the repository:
\begin{verbatim}
cc -std=c11 -O2 -Wall -Wextra -o /tmp/test_set_theory \
    tests/set_theory/test_set_theory.c && /tmp/test_set_theory
\end{verbatim}

\item \textbf{Expected result.} One \texttt{PASS} line for the file: none of the statements is proved.
\end{enumerate}
\end{document}
